\begin{enumerate}

    \item Validar prospectivamente el modelo con leads nuevos antes de adoptar sus resultados como criterio permanente de priorización. La evaluación realizada en este trabajo utilizó información histórica y demostró, mediante 1.6 millones de simulaciones de Monte Carlo, que el ordenamiento del modelo supera la selección aleatoria en todos los niveles de capacidad analizados. Sin embargo, una prueba prospectiva permitiría determinar si esta ventaja se mantiene cuando la priorización se utiliza dentro de la operación comercial. Se recomienda comparar durante un periodo definido los resultados obtenidos con el modelo frente a la gestión sin un mecanismo sistemático de priorización, utilizando indicadores como \textit{lift}, conversiones concentradas y cantidad de leads atendidos por conversión.

    \item Evaluar la capacidad de generalización del modelo con información de nuevos proyectos residenciales y, posteriormente, en otros contextos comerciales. El estudio se desarrolló con información histórica de seis proyectos de una misma empresa, por lo que la aplicabilidad de los patrones identificados fuera de esta población debe comprobarse empíricamente. Esta validación permitiría determinar qué relaciones se mantienen en nuevos desarrollos y cuáles dependen de las características particulares de los proyectos, las estrategias de captación o la población analizada.

    \item Mejorar la trazabilidad de la información que se incorpora durante el recorrido del lead por el proceso comercial. El análisis mostró que varios campos se completan de manera progresiva y que su ausencia puede contener información relacionada con el avance del prospecto. Se recomienda registrar, cuando sea posible, el momento en que cada dato se incorpora al CRM y distinguir entre información aún no recopilada y campos que permanecen vacíos por problemas de registro. Esto permitiría estudiar con mayor precisión el efecto de la completitud y evitar que ambos tipos de ausencia sean tratados como una misma condición en modelos posteriores.

    \item Ampliar el registro estructurado de las interacciones tempranas con los leads. Variables como la cantidad de intentos de contacto, el momento de cada interacción, el canal utilizado y el resultado obtenido podrían aportar información adicional antes de la etapa de priorización. Su incorporación en investigaciones posteriores permitiría evaluar si el comportamiento inicial del prospecto mejora la capacidad predictiva respecto de las variables disponibles en el presente estudio, siempre que estas señales puedan conocerse antes del desenlace comercial.

    \item Monitorear el desempeño del modelo conforme se incorporen nuevos datos y cambien las condiciones de captación y comercialización. La información analizada abarca distintos proyectos y periodos, por lo que las relaciones aprendidas por el modelo pueden modificarse ante cambios en campañas, canales, perfiles de clientes o características del mercado. Se recomienda dar seguimiento a métricas como PR-AUC, Brier score y \textit{lift} en los niveles de capacidad utilizados por el equipo comercial, y recalibrar o volver a entrenar el modelo cuando se observen cambios sostenidos en su desempeño.

    \item Evaluar en estudios posteriores la utilidad comercial de los segmentos obtenidos mediante K-Means. La segmentación permitió identificar seis perfiles de leads a partir de sus características, pero su construcción fue descriptiva y no demuestra que estrategias comerciales diferenciadas produzcan mejores resultados para cada grupo. Se recomienda utilizar estos segmentos como base para diseñar y evaluar pruebas controladas de mensajes, contenidos o formas de seguimiento, con el fin de determinar si la caracterización obtenida puede traducirse en diferencias observables en la respuesta comercial.

\end{enumerate}